\section{\texttt{vmm\_unmap\_page} 与 \texttt{vmm\_unmap\_identity}}

\subsection{解除单页映射}

\codefile{src/kernel/mm/vmm.c}
\begin{lstlisting}[style=cstyle]
void vmm_unmap_page(uint32_t virt) {
    if ((virt & 0xFFFu)) return;           // 对齐检查

    uint32_t pdi = pd_index(virt);
    uint32_t pde = g_kernel_pd.entries[pdi];
    if (!(pde & PDE_PRESENT)) return;      // 页表不存在

    uint32_t pt_phys = pde & ~0xFFFu;
    page_table_t* pt = (page_table_t*)PHYS_TO_VIRT(pt_phys);
    pt->entries[pt_index(virt)] = 0;       // PTE 置 0

    if (g_vmm_ready) vmm_invlpg(virt);    // 刷新 TLB
}
\end{lstlisting}

\begin{warnbox}
\texttt{vmm\_unmap\_page} 只断开虚拟→物理的映射，\textbf{不释放物理页}。
物理页可能被多处引用（如共享内存），释放是调用者的责任。
典型模式：先 \texttt{vmm\_get\_physical()} 拿到物理地址，再 \texttt{vmm\_unmap\_page()}，
最后 \texttt{pmm\_free\_page()} 归还物理页。
\end{warnbox}

\subsection{拆除 Identity Mapping}

\codefile{src/kernel/mm/vmm.c}
\begin{lstlisting}[style=cstyle]
void vmm_unmap_identity(void) {
    if (!g_vmm_ready) return;

    uint32_t first_kernel_pde = pd_index(KERNEL_VIRT_OFFSET); // 768
    uint32_t cleared = 0;

    for (uint32_t i = 0; i < first_kernel_pde; i++) {
        if (g_kernel_pd.entries[i] & PDE_PRESENT) {
            g_kernel_pd.entries[i] = 0;
            cleared++;
        }
    }

    // 重新加载 CR3 做全局 TLB 刷新
    uint32_t pd_phys = VIRT_TO_PHYS(
        (uint32_t)(uintptr_t)&g_kernel_pd);
    vmm_load_page_directory(pd_phys);
}
\end{lstlisting}

\begin{cpustate}
\texttt{vmm\_unmap\_identity()} 执行后：
\begin{itemize}
  \item PDE[0]--PDE[767] 全部清零
  \item 虚拟地址 \hex{00000000}--\hex{BFFFFFFF} 访问将触发 \#PF
  \item CR3 被重新加载，TLB 全局刷新
\end{itemize}
\end{cpustate}

% ============================================================================

